\section{Normal-measure Prikry constructions}
\label{annex:p06-prikry-forcing-cofinality}

All combinatorial constructions in this section are carried out inside the
ground \(M\) of \cref{def:p06-prikry-ground}. No ambient completeness of
the ground ultrafilter is used. The recursive forcing density and truth
facts are those of \cref{thm:p06-ftr-density,thm:p06-ftr-truth,thm:p06-ftr-zfc}.

\begin{annexitem}{annex:p06.prikry.measure-tools}{thm:p06-prikry-measure-tools}{Completeness, regularity and diagonal tails}
\begin{proof}
For any \(X\subseteq\kappa\) of size less than \(\kappa\), intersect
the complements of its singletons. Nonprincipality and completeness give
\(\kappa\setminus X\in U\), so \(X\notin U\). Consequently every
measure-one set has size \(\kappa\) and is unbounded, and every tail
above an ordinal belongs to \(U\).
If \(\operatorname{cf}(\kappa)=\lambda<\kappa\), a cofinal sequence
partitions \(\kappa\) into \(\lambda\) bounded sets. Their complements
are all in \(U\), whose intersection would be empty, a contradiction.
Thus \(\kappa\) is regular.

If \(2^\lambda\geq\kappa\), \(\lambda<\kappa\), choose an injection
\(e:\kappa\to\mathcal P(\lambda)\). For each \(i<\lambda\), one of
\(\{\alpha:i\in e(\alpha)\}\) and its complement lies in \(U\).
Intersect these fewer than \(\kappa\) chosen sets. All \(e(\alpha)\)
on the intersection have the same membership bits, so injectivity says
the intersection has at most one point, contradicting its membership
in \(U\). Hence \(\kappa\) is a strong limit. In particular, above
each ordinal below \(\kappa\) there is a ground cardinal still below
\(\kappa\).

If the complement of \(\triangle_{\xi<\kappa}A_\xi\) were in \(U\),
remove zero and send \(\beta\) to the least \(\xi<\beta\) for which
\(\beta\notin A_\xi\). Normality makes this regressive map constant
on a measure-one set. That set would miss one fixed \(A_\xi\in U\),
contradicting the filter property. This proves diagonal closure.

For finite stems, put
\(E_\xi=\bigcap\{B_t:t\in[\xi+1]^{<\omega},\ t\ne\varnothing\}\).
There are fewer than \(\kappa\) such finite stems, so \(E_\xi\in U\).
Take \(D=B_{\varnothing}\cap\triangle_{\xi<\kappa}E_\xi\).
If \(\max(t)<\beta\in D\), then
\(\beta\in E_{\max(t)}\subseteq B_t\). This is the asserted finite-stem
tail diagonalization with its index bound made explicit.

Conversely, finite-stem tail thinning implies diagonal closure: for a
given \(A_\xi\in U\), assign \(B_{\{\xi\}}=A_\xi\), with arbitrary
measure-one sets at the other stems. The resulting \(D\) lies in every
\(A_\xi\) whose index is below a point of \(D\), so
\(D\subseteq\triangle_{\xi<\kappa}A_\xi\).
If diagonal closure holds and a regressive \(f\) on \(A\in U\) has
no measure-one fiber, each complement
\(B_\xi=\kappa\setminus\{\alpha\in A:f(\alpha)=\xi\}\) belongs
to \(U\). A nonzero \(\alpha\) in
\(A\cap\triangle_{\xi<\kappa}B_\xi\) must both belong and not
belong to \(B_{f(\alpha)}\), because \(f(\alpha)<\alpha\).
This contradiction proves the regressive clause, completing all three
directions of equivalence.
\end{proof}
\end{annexitem}

\begin{annexitem}{annex:p06.prikry.homogeneity}{thm:p06-prikry-homogeneity}{Induction on finite arity}
\begin{proof}
First fix \(n\). For \(n=1\), partition \(A\) into the fewer than
\(\kappa\) color fibers. One fiber belongs to \(U\): if none did,
the intersection of their complements would be measure one and miss
\(A\). For \(n=0\), the singleton domain consisting of the empty
tuple is already homogeneous.

Suppose the result holds for \(n\geq1\) and let
\(c:[A]^{n+1}\to\gamma\). For each \(t\in[A]^n\), partition the
measure-one tail of \(A\) above \(\max(t)\) by the value
\(c(t\mathbin{\frown}\langle\beta\rangle)\). Choose a color \(i_t\)
and \(B_t\in U\) consisting of that color's tail points. Finite-stem
diagonalization gives \(D\in U\) with
\(\beta\in D\), \(\beta>\max(t)\) implying \(\beta\in B_t\).
Apply the induction hypothesis to \(t\mapsto i_t\) on \([A\cap D]^n\)
and choose homogeneous \(H\in U\), \(H\subseteq A\cap D\).
Every increasing \((n+1)\)-tuple in \(H\) has its last point in the
chosen tail of its first \(n\) points, so its color is the common
\(i_t\). This proves the induction step.

For the original finite-sequence coloring, choose a homogeneous
\(H_n\in U\) for each arity and take \(H=\bigcap_{n<\omega}H_n\).
Countable completeness, since \(\omega<\kappa\), keeps \(H\) in
\(U\). Restriction preserves each arity's homogeneity.
\end{proof}
\end{annexitem}

\begin{annexitem}{annex:p06.prikry.order-bridge}{thm:p06-prikry-order-bridge}{Increasing enumeration preserves the actual forcing structure}
\begin{proof}
Every finite subset of an ordinal has a unique increasing enumeration;
the inverse map takes the range of the sequence. Both maps are ground
functions. The condition on the upper part is unchanged: every member of
\(A\) is above the finite set exactly when it is above its enumeration.
For finite sets \(p,q\), end extension of \(q\) by \(p\) says that
the enumeration of \(q\) is an initial segment of that of \(p\).
The added sequence entries are exactly \(p\setminus q\). Consequently
each of the three native full-extension clauses is equivalent to the
corresponding sequence-stem clause. Equality of the first coordinates
and inclusion of upper parts likewise give both direct-extension clauses.
This proves preservation and reflection of stronger-than in both orders;
it does not equate the two displayed \(\leq\) conventions.

A common stronger condition transfers by the bijection, so compatibility
is preserved and reflected. A set is dense when every condition has a
stronger member of that set. The bijection carries each such witness to
the other presentation, proving preservation of density. A filter is
closed toward weaker conditions and directed toward stronger conditions;
both properties transfer by the same two implications. Because the
bijection and its inverse belong to \(M\), ground dense sets also transfer
to ground dense sets. Meeting every such set is therefore equivalent,
which proves the generic-filter assertion in both directions.

Translate a native name recursively by replacing every supporting
condition by its sequence-stem image and recursively translating each
member name. Well-foundedness of name rank gives a unique translation;
the inverse recursion undoes it by induction on rank. The atomic
recursive forcing clauses quantify over stronger conditions and the
support pairs of names. Their quantified conditions correspond
bijectively, and the lower-rank member names correspond by the induction
hypothesis, so the atomic membership and equality clauses agree. Induction
on formulas then carries this equivalence through the connective,
negation and quantifier clauses: each stronger-condition witness and each
name witness has a corresponding image and inverse. Finally, induction
on name rank in the evaluation recursion shows that translated names
have equal evaluations under corresponding filters. Thus the same
forcing conclusions, not merely the same cardinality of condition sets,
are obtained in the two presentations.
\end{proof}
\end{annexitem}

\begin{annexitem}{annex:p06.prikry.orders}{thm:p06-prikry-orders}{Partial orders, direct bounds, and compatible stems}
\begin{proof}
Reflexivity is immediate. For transitivity, if \(r\leq q\leq p\),
the stem of \(r\) end-extends that of \(p\); its new points are either
new points of \(q\), hence in \(A_p\), or added after \(q\), hence
in \(A_q\subseteq A_p\). Its tail is contained in \(A_p\).
Antisymmetry follows from mutual end extension and mutual tail inclusion.
The direct order has the same checks with a fixed stem and is contained
in the full order.

In a nonempty direct descending sequence of length \(\lambda<\kappa\),
all stems equal \(s\), and the intersection of the \(\lambda\) tails
belongs to \(U\), stays above \(s\), and gives a direct lower bound.
The empty sequence has the greatest condition as a lower bound.
For any fixed stem, finitely many conditions have common lower bound
with their tails intersected. There are exactly \(\kappa\) possible
finite stems. An antichain can have at most one condition of each stem,
so it has size at most \(\kappa\); this is \(\kappa^+\)-cc.

Choose any strictly increasing sequence \(\alpha_n<\kappa\) in the
ground and let \(p_n\) have stem
\(\langle\alpha_0,\ldots,\alpha_{n-1}\rangle\) and tail above its
last point. Then \(p_{n+1}\leq p_n\), but a lower bound would have
a finite stem extending every one of these unbounded-length stems.
That is impossible. This full-order failure is consistent with direct
closure because the sequence is not direct descending.
\end{proof}
\end{annexitem}

\begin{annexitem}{annex:p06.prikry.sequence}{thm:p06-prikry-sequence}{Dense requirements, cofinality and filter recovery}
\begin{proof}
Any two stems occurring in \(G\) have a common end extension, by
directedness; hence one is an initial segment of the other. Their union
is an increasing sequence with domain contained in \(\omega\).
It is an actual extension set, with canonical ground name
\(\dot C=\{(\check{\langle n,\alpha\rangle},p):p\in P_U,
 n<\operatorname{length}(s_p),\ s_p(n)=\alpha\}\)
for its graph. Evaluation gives exactly the union of the stem graphs.
For each \(n\), the conditions whose stems have length at least \(n\)
are dense: append enough successive points from the unbounded tail and
then trim the tail above the last point. Genericity gives domain exactly
\(\omega\). For each \(\alpha<\kappa\), the conditions whose stem has
a point above \(\alpha\) are dense by the same extension argument.
Thus \(C_G\) is cofinal. Since \(\kappa\) is regular uncountable in
the ground, such a sequence cannot belong to \(M\); in the extension
its cofinality is exactly \(\omega\), because no limit ordinal has
a finite cofinal subset.

For \(A\in U\), the conditions with tail contained in \(A\) are
dense: intersect the old tail with \(A\) and keep its stem. A condition
\((s,B)\in G\) with \(B\subseteq A\) ensures that all later Prikry
points are in \(B\). Indeed, for any other stem occurring in \(G\),
a common strengthening with \((s,B)\) places all its points after
\(s\) in \(B\). Hence \(C_G\) is eventually in \(A\).

Every \((s,A)\in G\) therefore satisfies the displayed reconstruction
condition. Conversely suppose \(p=(s,A)\) satisfies it. It is compatible
with every \(r=(t,B)\in G\): both stems are initial segments of
\(C_G\); take the longer one, whose points after the shorter stem lie
in the shorter condition's tail, and intersect the tails above the longer
stem. This is a common strengthening. The set
\(\{q:q\leq p\text{ or }q\perp p\}\) is dense and belongs to \(M\).
Its intersection with \(G\) cannot be an incompatible condition, so
some \(q\in G\) strengthens \(p\). Upward closure puts \(p\) in \(G\).
This proves equality with the reconstructed filter.

The sequence is definable from \(G\), and the displayed Separation
definition recovers \(G\) from \(C_G\) and the ground poset. Any
transitive \(\mathsf{ZFC}\) model containing the ground and \(C_G\)
therefore contains \(G\), and then all evaluations of ground names by
well-founded recursion. Thus the generated extension \(M[C_G]\) is
exactly \(M[G]\), as asserted.
\end{proof}
\end{annexitem}

\begin{annexitem}{annex:p06.prikry.lemma}{thm:p06-prikry-lemma}{Decision colors, diagonal tails and density}
\begin{proof}
Fix \(p=(s,A)\) and a formula with ground names. For each
\(t\in[A]^{<\omega}\), assign color zero if some condition with exact
stem \(s\mathbin{\frown}t\) and tail contained in \(A\) forces
\(\varphi\); assign color one if such a condition forces
\(\neg\varphi\); assign color two if neither exists. The first two
possibilities cannot both occur: exact-stem conditions are compatible,
and monotonicity and consistency of recursive forcing forbid opposite
decisions in compatible conditions. For colors zero and one choose a
witnessing tail \(B_t\in U\); for color two choose the tail of \(A\)
above \(t\). In every case \(B_t\subseteq A\) and lies above the new
stem. These are ground choices: fixed-formula forcing is definable in
\(M\), and the family of finite stems is a set there.

Homogenize this three-color map at every finite arity to get
\(H_0\in U\), \(H_0\subseteq A\). Apply finite-stem diagonalization
to all the \(B_t\), including the empty stem, and intersect the resulting
set with \(H_0\); call it \(H\). Thus whenever \(t\subseteq H\)
is nonempty,
\[
 H\setminus(\max(t)+1)\subseteq B_t,
 \qquad H\subseteq B_{\varnothing}.
\]
Put \(q=(s,H)\leq^*p\). By the density of decisions in the full
forcing order, choose \(r\leq q\) deciding the formula. Its stem is
\(s\mathbin{\frown}t\) for some \(t\in[H]^n\), and its tail is
contained in \(A\). The homogeneous color at arity \(n\) is therefore
zero or one, say \(i\).

For every \(u\in[H]^n\), the pure-tail condition
\((s\mathbin{\frown}u,H\setminus(\max(u)+1))\) strengthens the
selected witness with tail \(B_u\), and thus forces the same decision
\(\varphi\) if \(i=0\), or \(\neg\varphi\) if \(i=1\).
For \(n=0\), use \(H\subseteq B_{\varnothing}\).
Any extension \(r'\leq q\) with at least \(n\) added stem points
strengthens the witness for its first \(n\) new points: all subsequent
stem points and its tail lie in the corresponding \(B_u\) by diagonal
thinning. An extension with fewer added points can append more from its
measure-one tail until it has \(n\) of them. Conditions forcing decision
\(i\) are consequently dense below \(q\). Density closure of the
recursive forcing relation, \cref{thm:p06-ftr-density}, says that
\(q\) itself forces that decision.

The argument constructs the homogeneous and diagonal sets before using
density. It does not assume that a merely compatible exact-stem condition
already has the same forcing value on all its full extensions, and it
does not appeal to a truth lemma to define the decision colors.
\end{proof}
\end{annexitem}

\begin{annexitem}{annex:p06.prikry.strong-dense}{thm:p06-prikry-strong-dense}{Uniform entry into a dense open set}
\begin{proof}
For \(t\in[A]^{<\omega}\), give color one if some exact-stem
\((s\mathbin{\frown}t,B_t)\), with \(B_t\subseteq A\), belongs
to \(D\), and color zero otherwise. Choose such a tail for color one,
and the tail of \(A\) for color zero. Homogenize every arity and
diagonalize all chosen tails as in the preceding construction, obtaining
\(H\in U\), \(H\subseteq A\). This uses only normality tools, not
the conclusion of the decision lemma.

Density gives \(r\in D\) below \((s,H)\); let \(n\) be its number
of added stem points. The homogeneous color at arity \(n\) is one.
For every \(t\in[H]^n\), the condition with that stem and pure tail
\(H\setminus(\max(t)+1)\) strengthens the chosen \(D\)-witness
and belongs to \(D\), since \(D\) is downward open. If \(m>n\),
the analogous condition for a longer \(u\in[H]^m\) strengthens the
one for its first \(n\) points, and also belongs to \(D\). When
\(n=0\), the empty-stem tail diagonalization supplies \((s,H)\in D\)
and all its further pure-tail stem additions. This proves the stated
uniform finite-level property.
\end{proof}
\end{annexitem}

\begin{annexitem}{annex:p06.prikry.no-bounded}{thm:p06-prikry-no-bounded}{A fixed stem decides a bounded subset name}
\begin{proof}
Fix \(p=(s,A)\) forcing \(\dot X\subseteq\check\lambda\), where
\(\lambda<\kappa\). Recursively for \(\alpha<\lambda\), take a
direct lower bound of the previous conditions, keeping \(s\), then
apply the proved Prikry lemma to decide
\(\check\alpha\in\dot X\). At every limit stage the lower bound
is the same stem with the intersection of the previous tails; there
are fewer than \(\kappa\) tails. At the end intersect all the chosen
tails and the original \(A\), obtaining \(q\leq^*p\).
Let \(Y=\{\alpha<\lambda:q\forces\check\alpha\in\dot X\}\).
This set belongs to \(M\) by Separation using the fixed-formula
definability of recursive forcing, and each omitted \(\alpha\) is
forced not to belong to \(\dot X\).

We justify the equality conclusion rather than inferring a forcing
statement from an unqualified universal claim about generics. If a
strengthening of \(q\) forced \(\dot X\ne\check Y\), Extensionality
and the forced boundedness would give, on a further strengthening, a
name for a point in their symmetric difference and in \(\check\lambda\).
The existential forcing clause supplies such witness conditions densely.
The atomic membership clause for \(\check\lambda\) then supplies,
on another strengthening, equality of that name with some
\(\check\alpha\), \(\alpha<\lambda\): the check name has precisely
these ground member names with unit support. This contradicts the
already fixed membership decision for that \(\alpha\). No strengthening
forces inequality; the dense decision set and classical recursive forcing
therefore give \(q\forces\dot X=\check Y\).
For \(\lambda=0\), forced inclusion in the empty set already gives
this equality without any recursion.

The same construction works below every strengthening of \(p\), so
conditions deciding \(\dot X\) as a ground subset form a ground dense
set below \(p\). A generic containing \(p\) meets it, using its union
with the conditions incompatible with \(p\) as a global dense set.
By the truth lemma, every subset of \(\lambda\) in \(M[G]\) is
therefore in \(M\). The converse inclusion is automatic. In particular
\(\lambda=\omega\) gives no new reals. Any bounded subset of \(\kappa\)
has a bound \(\lambda<\kappa\), an ordinal already in the ground,
so this proves the full bounded-subset assertion.
\end{proof}
\end{annexitem}

\begin{annexitem}{annex:p06.prikry.preservation}{thm:p06-prikry-preservation}{Graph codes and possible-value antichains}
\begin{proof}
First let \(\lambda<\kappa\) be a ground cardinal. A hypothetical
extension surjection \(f:\gamma\to\lambda\), \(\gamma<\lambda\),
has graph included in \(\gamma\times\lambda\). In the ground this
product injects into some ordinal \(\theta<\kappa\). The graph's
image is a bounded subset of \(\kappa\), hence is in the ground by
the preceding theorem. Decoding by the ground injection puts \(f\)
itself in \(M\), contradicting that \(\lambda\) was a cardinal there.

At \(\lambda=\kappa\), a surjection from \(\gamma<\kappa\) would
also give a surjection onto any ground cardinal
\(\mu\) with \(|\gamma|^M<\mu<\kappa\): keep values below \(\mu\)
and replace larger values by zero. Such \(\mu\) exists by the ground
strong-limit property proved earlier. This new function has a bounded
graph code and would be a ground surjection, again impossible. This is
why the no-bounded-subsets argument at \(\kappa\) uses its ground
limit-cardinal property, not the mistaken claim that the original
unbounded graph is itself bounded.

For a ground cardinal \(\lambda>\kappa\), suppose some \(p\) forces
that \(\dot f:\check\gamma\to\check\lambda\) is onto, with
\(\gamma<\lambda\). For each \(\xi<\gamma\), choose in \(M\) a
maximal antichain below \(p\) deciding \(\dot f(\check\xi)\) as a
ground ordinal. Such decisions are dense by the check-name membership
clause; maximal antichains exist by Choice. The chain condition bounds
each antichain by \(\kappa\), so the ground set \(S_\xi\) of possible
values has size at most \(\kappa\). Maximality makes conditions forcing
membership in \(\check S_\xi\) dense below \(p\); hence \(p\) forces
that membership. Thus \(p\) forces the range inside
\(T=\bigcup_{\xi<\gamma}S_\xi\), whose ground cardinality is at most
\(\max\{|\gamma|^M,\kappa\}<\lambda\). A ground point of
\(\lambda\setminus T\) cannot be in the range, contradicting surjectivity.
This preserves every cardinal above \(\kappa\), including singular ones.
Every ground noncardinal ordinal has a ground bijection with a smaller
ordinal, which persists; since forcing adds no ordinals, there are no
new cardinal ordinals either.

For the cofinality statement, fix a nonzero ground limit \(\delta\)
and put \(\mu=\operatorname{cf}^M(\delta)\). A ground increasing
cofinal map \(h:\mu\to\delta\) persists, so cofinality can only
decrease. If \(\mu<\kappa\) and an extension cofinal map
\(f:\nu\to\delta\), \(\nu<\mu\), existed, send each \(f(\xi)\)
to the least \(\eta<\mu\) for which \(f(\xi)<h(\eta)\).
This gives a cofinal map \(g:\nu\to\mu\). Its graph has a bounded
ground code below \(\kappa\); no new bounded subsets puts \(g\) in
\(M\), contradicting regularity of the ground cofinality \(\mu\).

If \(\mu>\kappa\), repeat the possible-value antichain construction
for a proposed cofinal map \(\dot f:\nu\to\delta\), \(\nu<\mu\).
The union of its possible values has ground size at most
\(\max\{|\nu|^M,\kappa\}<\mu\), and hence is bounded in \(\delta\)
by the definition of ground cofinality. The condition would force a
cofinal range inside a fixed ground bound, a contradiction. Finally if
\(\mu=\kappa\), the map \(n\mapsto h(C_G(n))\) is cofinal in
\(\delta\), so its extension cofinality is \(\omega\). No finite set
is cofinal in a limit ordinal. Zero and successor cofinalities are
unchanged by their ordinal definitions. This completes the classification.
\end{proof}
\end{annexitem}

\begin{annexitem}{annex:p06.prikry.measure-loss}{thm:p06-prikry-measure-loss}{The new family lies outside the ground completeness quantifier}
\begin{proof}
For each ground cardinal \(\lambda<\kappa\), its power set is unchanged,
and the cardinality of that ground power set is unchanged by cardinal
preservation. The ground strong-limit inequality
\(2^\lambda<\kappa\) therefore persists for every extension cardinal
\(\lambda<\kappa\); there are no new cardinal ordinals to check.
The cardinal \(\kappa\) is still uncountable but now has cofinality
\(\omega\). The regularity argument for a nonprincipal
\(\kappa\)-complete ultrafilter applies in any \(\mathsf{ZFC}\) model,
so no such ultrafilter can witness measurability of this singular cardinal
in \(M[G]\).

For every \(n\), the tail \(A_n=\kappa\setminus(C_G(n)+1)\) belongs
to the old \(U\): its endpoint is a ground ordinal and all ground tails
are in \(U\). Cofinality gives \(\bigcap_nA_n=\varnothing\).
The indexed family cannot belong to \(M\), since its empty intersection
would contradict ground countable completeness. Hence the old measure's
completeness is not absolute to the new families of the extension.
\end{proof}
\end{annexitem}

\begin{annexitem}{annex:p06.prikry.mathias}{thm:p06-prikry-mathias}{Eventual containment meets every ground dense set}
\begin{proof}
Every ground tail of \(\kappa\) lies in \(U\); eventual containment
in these tails makes \(C\) cofinal. Define \(G(C)\) by the displayed
stem/tail condition. It contains the greatest condition. Given two
members, use their longer initial stem and the intersection of their
tails above it; the intervening points of \(C\) lie in the shorter
member's tail. This common strengthening is still in \(G(C)\).
If \(q\in G(C)\) strengthens \(p\), then \(p\)'s shorter stem
is initial in \(C\); the finitely many intervening points are in its
tail by the order definition, and all subsequent points are in that
tail because the tail of \(q\) is contained in it. Thus \(p\in G(C)\).
So \(G(C)\) is a filter.

Fix a ground dense downward-open \(D\). For each finite increasing
stem \(s\), apply the proved strong dense-set lemma to the condition
with that stem and the entire ground tail above it. Choose
\(A_s\in U\) above \(s\), and \(n_s<\omega\), so that every
stem addition from \(A_s\) of length at least \(n_s\), with pure
\(A_s\)-tail, belongs to \(D\). These choices are a ground set.
Finite-stem diagonalization gives \(K\in U\) such that
\(\beta\in K\), \(\beta>\max(s)\) implies \(\beta\in A_s\),
with \(K\subseteq A_{\varnothing}\) as well.

Choose \(N\) with \(C(n)\in K\) for every \(n\geq N\), and put
\(s=C\mathbin{\upharpoonright}N\). Every finite initial segment is
in \(M\), since its ordinal entries are ground and finite tuples are
ground sets. All remaining points of \(C\) lie in \(A_s\), by the
diagonal condition (including the empty-stem convention). Let \(t\)
be the next \(n_s\) points. The strong lemma puts
\((s\mathbin{\frown}t,A_s\setminus(\max(t)+1))\) in \(D\), and
the rest of \(C\) lies in its tail, so it is also in \(G(C)\).
If \(n_s=0\), use \((s,A_s)\) instead. Hence \(G(C)\cap D\ne\varnothing\).

For an arbitrary ground dense set \(E\), its downward closure
\(\{q:(\exists p\in E)\ q\leq p\}\) is dense open. A member of
\(G(C)\) in that closure lies below a member of \(E\), which belongs
to the filter by upward closure. Thus \(G(C)\) meets every required
ground dense set and is \(M\)-generic.
\end{proof}
\end{annexitem}
